\section{The bracket, the Jones polynomial, and crossing number}
\label{sec:bracket}

The calculation in this section follows the state model of Louis H. Kauffman,
\emph{State models and the Jones polynomial}, Topology \textbf{26} (1987),
395--407, doi:10.1016/0040-9383(87)90009-7. All sums, degrees, and component
counts below are integer calculations on finite diagrams.

\subsection{Resolving crossings}

Take a nonempty diagram in the oriented plane, compactified by an exterior
face. At a crossing mark the two corners swept out on turning
counterclockwise from each outgoing overpassing ray to the adjacent
underpassing ray. The \emph{$A$-pairing} joins the two rays of each marked
corner. The other adjacent pairing is the $B$-pairing. In a downward braid
with its top-left strand passing over the top-right strand, the $A$-pairing
consists of two vertical strands. This convention fixes the handedness of all
the formulas below.

Number the $c$ crossings. A state $S$ is a $c$-bit mask, with a set bit
selecting $B$. Write $b(S)$ for its number of set bits and $h(S)$ for the
number of circles after making the selected pairings. To compute $h$, form
the graph on half-edges with one edge per diagram arc and one per local
pairing. Count connected components and add the number of crossing-free
circles. Each component of this graph is an alternating cycle of the two
edge types, allowing a cycle of length two. Thus each is one resolved
circle. With $\delta=-A^2-A^{-2}$ in $\mathbb Z[A,A^{-1}]$ put
\[
 B_D=\sum_{S<2^c} A^{c-2b(S)}\delta^{h(S)-1}.
\]
In particular $B_{\bigcirc}=1$. Partitioning masks by a chosen bit proves the
local identity $B_{\times}=A B_{A\text{-pairing}}+
A^{-1}B_{B\text{-pairing}}$. A fresh circle multiplies $B$ by $\delta$.
These formulas determine the sum on nonempty diagrams by induction on $c$.

\begin{theorem}[Regular isotopy and the Jones skein relation]
Let $w$ be the sum of the oriented crossing signs, positive for the
downward braid crossing just specified. Then
\[
 F_D(A)=(-A^3)^{-w(D)}B_D(A)
\]
is unchanged by all oriented Reidemeister moves. For
$t^{1/2}=A^{-2}$ it defines a Laurent polynomial $V_D(t)$ in $t^{1/2}$
with $V_{\bigcirc}=1$ and
\[
 t^{-1}V_{+}-t V_{-}=(t^{1/2}-t^{-1/2})V_{0},
\]
where $0$ is the orientation-preserving smoothing.
\end{theorem}
\begin{proof}
For a rectangular patch resolve its crossings but keep its boundary
endpoints. A resolved patch is a pairing of endpoints and a number of
internal circles. Give each internal circle weight $\delta$. Stacking
patches composes pairings by tracing paths and multiplies by $\delta$ for
each newly closed circle. This operation is associative since tracing in the
fully glued graph gives the same paths and closed components in either
order. Equality of two local expansions survives insertion into any
nonempty closed diagram: fix the outside state, glue each pairing to the
same outside arcs, and apply the common weight $\delta^{h-1}$ to the
resulting circles. There is at least one circle using the boundary when the
patch has endpoints.

For $k$ vertical strands write $I$ for the identity pairing and $e_i$ for a
cap and cup at positions $i,i+1$, with every other strand vertical.
Tracing the affected endpoints gives
\[
 e_i^2=\delta e_i,\qquad
 e_i e_{i+1}e_i=e_i,\qquad e_{i+1}e_ie_{i+1}=e_{i+1}.
\]
In the first stacking, the interior cup and cap form one closed circle. In
the second, the outside cap and cup stay at $i,i+1$ and the third strand
traverses both interior turns with no closed circle. Reflection gives the
third equality. Disjoint cap-cup pairs commute. The two crossing
expansions are $R_i=AI+A^{-1}e_i$ and $R_i^{-}=A^{-1}I+Ae_i$.
Their product in either order is
$I+(A^2+A^{-2}+\delta)e_i=I$, which proves move II. For $e=e_i$, $f=e_{i+1}$, put $U=AI+A^{-1}e$ and $W=AI+A^{-1}f$. Expanding the two sides of move III gives
\[
 UWU-WUW=(A+A^{-1}\delta+A^{-3})(e-f)=0 .
\]
These are the unoriented move templates. To see that they cover type III,
the three strands cross pairwise, so their six boundary ends occur in
order $x,y,z,x,y,z$ up to relabeling. Choose three consecutive ends as
inputs so that their strands, read around the disk in one of the two
directions, are in descending overcrossing height order. Such a triple
exists because cyclic orders and their reverses exhaust the six
permutations. The local pictures are then the two three-crossing braid
words above or their reflections. Reflection interchanges $A,A^{-1}$
and preserves the identity. Rotation, or bending an endpoint along the
boundary, merely glues the same pairing onto both sides. This also covers
rotated type II patches.

In a positive curl the $A$-pairing detaches a circle, giving factor
$A\delta+A^{-1}=-A^3$. The opposite curl gives
$A^{-1}\delta+A=-A^{-3}$. The oriented signs cancel in move II because
the same pair of strands crosses twice with opposite local order. In
move III the signs of the three corresponding pairs are unchanged. Move I
changes $w$ by the matching sign; hence the normalization cancels its
factor. Reversing all arrows fixes $w$ and $F$. Each exponent in $F$ is
even since $w\equiv c\pmod 2$ and $\delta$ has even exponents.

For the skein identity the positive crossing has the oriented smoothing
in its $A$-term. In the negative crossing the weights are interchanged.
Write $U$ for the bracket of the oriented smoothing and $H$ for that of
the other pairing. If $u$ is the outside writhe, the two normalized
brackets are $(-A^3)^{-u}(-A^{-3})(AU+A^{-1}H)$ and
$(-A^3)^{-u}(-A^{3})(A^{-1}U+AH)$. Multiplying the first by $A^4$
and subtracting $A^{-4}$ times the second cancels $H$ and leaves
$(A^{-2}-A^{2})(-A^3)^{-u}U$. This is the asserted identity.
\end{proof}

\begin{proposition}[Exponents and elementary operations]
If $D$ has $r$ link components, $V_D\in
t^{(r-1)/2}\mathbb Z[t,t^{-1}]$. Mirroring replaces $t$ by $t^{-1}$.
For a disjoint diagram union one has $F_{D\sqcup E}=\delta F_D F_E$.
For the oriented connected sum along two chosen arcs one has
$F_{D\#E}=F_D F_E$.
\end{proposition}
\begin{proof}
Consider the three possible pairings at one crossing after all other
crossings have been smoothed. The exterior arcs in the complementary disk
give one of the two adjacent pairings. Indeed, two disjoint arcs in a
disk cannot pair alternating boundary ends: closing one arc by either
boundary interval separates the other two ends, by the polygonal
crossing-parity test. Thus changing $A$ to $B$ splits or merges one
circle. The leading exponent of the corresponding state monomial changes
by $-2+2(\pm1)$, and the internal powers of $\delta$ have steps of four.
All exponents after normalization therefore belong to a single residue
class modulo four.

Use the state that smooths all crossings compatibly with orientation.
Let $p$ and $n$ count positive and negative crossings. Its bit count is
$n$. Starting with the original components, make these oriented smoothings
one by one. At each step the incoming ends previously led to two specific
outgoing ends, and the new prescription exchanges these successors.
A transposition of successors splits one directed cycle in two or joins two
directed cycles in one. Hence this state's circle count satisfies
$h\equiv r+c\pmod2$. Its normalized exponent modulo four is
$c-2n+2(h-1)-3w=2(h-1-w)\equiv 2(r-1)\pmod4$ since
$w=p-n\equiv c\pmod2$. This proves the claimed grading.

In the mirror all corners swap roles, masks complement, $w$ changes sign,
and $\delta$ is fixed by $A\mapsto A^{-1}$. For the union, states are
independent pairs and the circle counts add. For the sum, cutting the two
marked state circles opens each into a path; the joining arcs make a
single circle out of the pair. Thus circle counts add minus one. Crossing
weights and writhe add in both constructions.
\end{proof}

\begin{proposition}[Value at one and parity of the determinant]
For a diagram with $r$ components,
$F_D(1)=V_D(1)=(-2)^{r-1}$. In particular every Jones polynomial is
nonzero. For knots $V_D(-1)$ is an odd integer.
\end{proposition}
\begin{proof}
We prove $B_D(1)=(-1)^c(-2)^{r-1}$ by resolving one crossing.
Outside its four ends, follow the original strands straight through every
other crossing. This gives a pairing of the four ends and some closed
components. Of the three possible local pairings, the one agreeing with
the outside pairing yields two additional circles, and each of the other
two yields one. If $d$ denotes the smaller total circle count,
the three numbers $(-2)^{r_P-1}$, for the original, $A$, and $B$
pairings $P$, therefore sum to
$(-2)^d+2(-2)^{d-1}=0$. The two smoothing terms have $c-1$ crossings;
the bracket recursion at $A=1$ proves the induction step. The case $c=0$
is the disjoint-circle formula. Normalization multiplies by $(-1)^w$
with $w\equiv c\pmod2$. Finally, for a knot $V_D$ has integral exponents,
and $1$ and $-1$ agree modulo two in every Laurent monomial. Thus
$V_D(-1)\equiv V_D(1)=1\pmod2$.
\end{proof}

\subsection{Extreme coefficients and the first Tait conjecture}

Write $a=h(0)$, $d=h(2^c-1)$ for the all-$A$ and all-$B$ circle counts.
A diagram is $A$-adequate if at every crossing the two $A$-smoothing arcs
belong to distinct all-$A$ circles. Define $B$-adequacy likewise. The
\emph{shadow} forgets over/under and retains the crossing vertices. A
crossing-free circle is a separate connected shadow.

\begin{lemma}[Dual state bound]
For a connected shadow and any pair of complementary states $S,\bar S$,
$h(S)+h(\bar S)\le c+2$.
\end{lemma}
\begin{proof}
Take the $h(S)$ circles as vertices and one edge per crossing between the
two incident state arcs, allowing loops. This graph is connected: collapsing
the state circles and restoring the crossing sites preserves the connectivity
of the original shadow. Select a spanning tree by scanning the edges and
retaining an edge exactly when it joins distinct components. Changing the
pairings first at tree edges merges two circles each time. Indeed, by
induction each current tree component has exactly one circle, and the next
edge joins different components. After $h(S)-1$ changes one circle remains.
Each remaining change increases the count by at most one. Consequently
$h(\bar S)\le1+c-(h(S)-1)$.
\end{proof}

\begin{theorem}[Span bounds and adequacy]
For every nonempty diagram all exponents of $B_D$ lie between
$-c-2(d-1)$ and $c+2(a-1)$. Under $A$-adequacy the upper endpoint occurs
with coefficient $(-1)^{a-1}$; under $B$-adequacy the lower occurs with
coefficient $(-1)^{d-1}$. On any connected shadow the span of a nonzero
$F_D$ in the variable $A$ is at most $4c$.
\end{theorem}
\begin{proof}
A change from $A$ to $B$ reduces the prefactor degree by two and
increases the circle count by at most one. It therefore cannot increase
the upper bound for the state's degree. If $D$ is $A$-adequate, the
first such change merges two circles and drops the upper bound by four.
Every state other than all-$A$ is reached by changes of this kind. Only
all-$A$ contributes at $c+2(a-1)$ and it contributes the stated sign.
Mirror the argument for the lower extreme. Subtract the degree bounds
and use $a+d\le c+2$ to obtain span at most $4c$. Writhe normalization
translates degrees by a constant.
\end{proof}

For the next theorem take a connected sphere diagram with $c>0$.
In half-edge notation its cyclic orders and edge reversal trace the
facial boundary walks. For a sphere map, $c-2c+f=2$ where $f$ is the
number of face cycles. One can check this input condition by counting
cycles. Geometrically it is Euler's formula: a spanning tree in a
connected plane graph has one face, and inserting each remaining edge in
its face splits that face into two. Diagram arcs can be polygonal with
rational vertices. A crossing is \emph{nugatory} in the face test if two
opposite corners belong to the same face. Equivalently a simple closed
curve can pass through that crossing, approaching in the two opposite
sectors, with its remaining arc in the complement. In the face-polygon
description join the two corner occurrences by an interior simple arc,
then identify their endpoints at the crossing. Conversely, the interior
of any such arc belongs to a single face. A diagram is reduced if this
test fails at every crossing.

\begin{theorem}[Crossing minimality, with a spanning-tree certificate]
Let $D$ be connected and alternating, meaning successive visits on each
component alternate over and under. The class of $F_D$ modulo $A^4+1$
is a unit monomial times the positive number of spanning trees of a
checkerboard graph of $D$. In particular a diagram joined to $D$ by
Reidemeister moves has connected shadow. If $D$ is also reduced, its
$A$-span equals $4c$ and every such diagram has at least $c$ crossings.
\end{theorem}
\begin{proof}
Take first $c>0$. Shade the $A$-corners black and the $B$-corners white.
Along an arc an over endpoint is paired with an under endpoint; because
the two outgoing rays point oppositely along that arc, the black side
continues into the black side at the other end. Thus each facial walk
has a single color and each edge has one face on each side of distinct
colors. The all-$A$ circles trace precisely the black facial walks, and
the all-$B$ circles the white ones. Hence $a+d=f=c+2$. Put one vertex
in each black face and one edge at each crossing joining the two black
corners. This is the connected graph of the dual state bound starting
from all-$A$. The white graph is the corresponding graph starting from
all-$B$.

Reduce coefficients modulo $A^4+1$ in the Laurent polynomial ring.
Here $A$ is a unit and $\delta=0$, so only one-circle states survive.
Write $T$ for the set of changed black pairings, of size $j$. Every
state circle is contained in one component of the spanning black subgraph
with edges $T$: away from changed crossings follow the original black
circle, and at a changed crossing one can pass only between its two
incident black vertices. Each vertex supplies some arc of the diagram
and is thus represented in a state circle. Consequently a one-circle
state implies connectivity of this subgraph, so $j\ge a-1$. Starting
instead from white circles gives $c-j\ge d-1$. As $a+d=c+2$, both
inequalities are equalities. A connected graph on $a$ vertices with
$a-1$ edges is a tree, by the edge-scan construction above. Conversely
every spanning tree gives exactly one circle by the same merging argument.
The surviving sum is therefore $\tau A^{c-2(a-1)}$ with $\tau\ge1$ the
spanning-tree count, and normalization multiplies it by a unit monomial.
This class is nonzero: reducing a monomial to a signed power of degree
$0,1,2,$ or $3$ gives unique coordinates since $A^4+1$ is monic and
$A^4=-1$. If a shadow has two or more components, partition them into two
nonempty diagrams. The disjoint-union calculation gives a factor
$\delta$, even if one component is drawn inside another. Its class is
zero. This proves the connectivity assertion for any competitor.

If $D$ is reduced, both black corners are in distinct faces and both white
corners are in distinct faces at every crossing. Thus $D$ is adequate on
both sides. Its extreme degrees give span $2c+2(a+d)-4=4c$.
Any certified equivalent diagram $E$ has this span and, as just proved,
connected shadow. The span bound gives $4c\le4c(E)$. For $c=0$ the
connected diagram is a circle, $F=1$; with a one-vertex, zero-edge
checkerboard graph all conclusions hold directly.
\end{proof}

The same extreme degrees give a chirality test. For a reduced connected
alternating diagram, equivalence to its oriented mirror implies
$3w=a-d$: the greatest and least exponents of $F$ must sum to zero.
When $c>0$, reducedness forces $a,d\ge2$, so $|a-d|\le c-2$.
In particular $3|w|\ge c$ excludes equivalence to the oriented mirror.

\begin{proposition}[Tutte specialization and signs for alternating links]
In the preceding checkerboard graph $G$, put $v=a$ and $e=c$. Write
$k(U)$ for the number of components of the spanning subgraph with edge
set $U$. Define its Tutte polynomial, for this connected $G$, by
\[
 T_G(X,Y)=\sum_{U\subseteq E}
 (X-1)^{k(U)-1}(Y-1)^{|U|-v+k(U)}.
\]
Then
\[
 B_D(A)=A^{e-2(v-1)}T_G(-A^4,-A^{-4}).
\]
After an overall sign and monomial shift, the coefficients of $F_D$ at
$A^{4j}$ have signs $(-1)^j$ wherever they are nonzero.
For an alternating knot $|V_D(-1)|$ counts the spanning trees of $G$,
so that count is odd.
\end{proposition}
\begin{proof}
Draw a small disk for each black face and thin bands for the edges in
$U$, ending at its cyclically marked corners. The boundary circuits use
the $B$-pairing at selected bands and the $A$-pairing elsewhere: a
band locally replaces two separate boundary arcs by the other smoothing.
For one connected component of this spanning subgraph, a neighborhood of
a spanning tree is a disk, by adding a new disk and its joining band at
each step. Each remaining edge lies in a face of the connected partial
plane graph. That face is a disk and its boundary in the thickened graph
is one circle; the new band cuts the face into two and adds one boundary
circle. (One can make this edge insertion with that component alone in
the sphere.) Thus a component with $v'$ vertices and $e'$ edges has
$2-v'+e'$ boundary circles, also valid for an isolated vertex. Summing
gives $h(U)=2k(U)-v+|U|$. For $i=k(U)-1$ and
$j=|U|-v+k(U)$ we have $h(U)-1=i+j$ and $|U|=v-1-i+j$.
Consequently the state term, divided by $A^{e-2(v-1)}$, is
$(A^2\delta)^i(A^{-2}\delta)^j$. Since
$1+A^2\delta=-A^4$ and $1+A^{-2}\delta=-A^{-4}$ this proves the formula.

For coefficient signs we need nonnegativity of the ordinary coefficients
of $T_G$. Its subset definition extends to disconnected graphs by changing
$k(U)-1$ to $k(U)-k(E)$. A loop $s$ contributes the factor
$1+(Y-1)=Y$ times $T_{G\setminus s}$, by pairing subsets with and
without $s$. For a bridge $s$, pair them instead with subsets in $G/s$:
omitting $s$ adds one free component relative to including it, so the
factor is $(X-1)+1=X$. For an edge of neither kind the two classes of
subsets give $T_{G\setminus s}+T_{G/s}$, as deletion preserves the total
component count, while contracting a selected non-loop decreases both
its vertex and selected-edge counts by one. The graph with no edges
has $T=1$. Induction now proves nonnegativity. Every monomial $X^iY^j$
contributing to degree $4(i-j)$ in the specialization has sign
$(-1)^{i+j}=(-1)^{i-j}$, proving the assertion. For a knot the exponent
grading places all powers of $F$ in residue zero modulo four.
Its class at $A^4=-1$ is exactly $V(-1)$; the spanning-tree evaluation
and parity assertion above conclude the proof.
\end{proof}

\subsection{Concrete knots and an infinite family}

Let $T_m$ be the closure of the downward two-strand braid $\sigma_1^m$,
$m\ge0$, the usual diagram of the $(2,m)$ torus link. Its strands
undergo $m$ transpositions, so there is one component for odd $m$ and
two for even $m$. Write $R=AI+A^{-1}e$. Since $Re=-A^{-3}e$,
induction, or telescoping $R^m-(AI)^m$, gives
\[
 R^m=A^m I+\sum_{j=0}^{m-1}(-1)^j A^{m-2-4j} e.
\]
Closing $I$ gives two circles of bracket $\delta$, while closing $e$
gives one of bracket $1$. For $m\ge1$ the $j=0$ term cancels part of
$A^m\delta$. Thus
\[
 F_{T_m}=(-A^3)^{-m}\left(-A^{m+2}
           +\sum_{j=1}^{m-1}(-1)^j A^{m-2-4j}\right).
\]
For $m\ge2$ the extreme exponents of $F$ are $2-2m$ and $2-6m$,
giving span $4m$ and chirality for the indicated orientation. The
polynomial modulo $A^4+1$ is nonzero by the spanning-tree result for this
alternating diagram. It follows that its crossing number in the
Reidemeister presentation is exactly $m$, including for even $m$.
The odd $m\ge3$ give pairwise distinct nontrivial knots of unbounded
crossing number. In particular the right trefoil has $V=t+t^3-t^4$,
its mirror has $V=t^{-1}+t^{-3}-t^{-4}$, and both require three crossings.

\begin{center}
\begin{tikzpicture}[line width=0.95pt,line cap=round]
\begin{scope}
\draw[ink] (-.85,0) .. controls (-.45,.87) and (.45,.87) .. (.85,0);
\draw[ink] (-.85,0)--(.85,0);
\draw[ink] (-.85,0) .. controls (-.45,-.87) and (.45,-.87) .. (.85,0);
\fill[sea] (-.85,0) circle (2.3pt) (.85,0) circle (2.3pt);
\node[font=\small] at (0,-1.01) {black graph};
\end{scope}
\begin{scope}[shift={(4.0,0)}]
\draw[ink] (0,.66)--(-.70,-.48)--(.70,-.48)--cycle;
\fill[sea] (0,.66) circle (2.3pt) (-.70,-.48) circle (2.3pt) (.70,-.48) circle (2.3pt);
\node[font=\small] at (0,-1.01) {white graph};
\end{scope}
\end{tikzpicture}
\end{center}
\noindent For the standard positive trefoil all-$A$ smoothing leaves two circles; all-$B$ leaves three. Its black graph therefore has three parallel edges, and its plane dual is a triangle. Choosing any one black edge (or omitting any one white edge) gives a spanning tree. This recovers $|V(-1)|=3$ and makes the extreme bracket exponents $5$ and $-7$ visible from the two circle counts.

Let $P=t+t^3-t^4$ and $P^*(t)=P(t^{-1})$. The granny knot, its mirror,
and the square knot have polynomials $P^2,(P^*)^2,PP^*$, respectively:
they are the sums of two right trefoils, two left trefoils, and one of
each. Their degree intervals $[2,8],[-8,-2],[-3,3]$ separate all three.
Each knot has crossing number six, since a six-crossing sum diagram
exists and its Jones span in $t$ is six.

For the figure-eight knot use the closed three-braid
$(\sigma_1\sigma_2^{-1})^2$, with zero writhe. Set
$x=A^{-2}$, $y=A^2$, $e=e_1$, $f=e_2$. The bracket expansion of one
factor is $Q=I+xe+yf+ef$. The normalized closure evaluation
$\operatorname{tr}$ sends $I$ to $\delta^2$, $e,f$ to $\delta$ and
$ef,fe$ to $1$, by direct tracing. It is cyclic since the seam of a
stacked closed rectangle can be taken before or after either factor.
Expanding $Q^2$ and using the cap-cup identities gives
\[
 \operatorname{tr}(Q^2)
  =(1+x^2+y^2)\delta^2+4(x+y)\delta+5
  =A^8-A^4+1-A^{-4}+A^{-8}.
\]
Thus $V=t^2-t+1-t^{-1}+t^{-2}$ and four crossings are necessary and
sufficient (the closure permutation is a three-cycle). Its mirror
equivalence also has a finite witness. If $s=\sigma_1$, $u=\sigma_2$,
then $\Delta=sus=usu$ conjugates $s$ to $u$ and $u$ to $s$, by that
same braid relation. Conjugation therefore changes $(su^{-1})^2$ to
$(us^{-1})^2$. Cyclically permuting the latter gives $(s^{-1}u)^2$,
the mirror. For closed braids cyclic permutation slides the marked seam
around the annulus, and conjugation reduces to cyclic permutation and
adjacent inverse cancellation. These observations supply the moves;
equality of the two Laurent polynomials alone would not supply them.

\paragraph{Arithmetic implementation.}
A table of $4c$ half-edges, plus $z$ isolated circles, suffices. Boolean
reachability closure by $4c$ rounds computes the state components and
their least representatives. Every state has $1\le h\le H=2c+z$:
each non-isolated component uses at least two half-edges.
Allocate exponent slots from $-4c-2H$ through $4c+2H$ to include
normalization. The sum of absolute coefficients before cancellation is
bounded by $2^c 2^{H}$, since one state has coefficient norm
$2^{h-1}$. Powers of $\delta$ can be accumulated by shifted addition.
Spanning-tree selection needs one edge scan with bounded component
recomputation; counting trees needs at most $2^c$ subsets. Each proof
inducts on a mask fold, a crossing-change list of length at most $c$, a
tree-edge scan, or a move list with its given length. For the evaluation at one use a postorder fold over the ternary table of partially resolved states of the original diagram. For deletion--contraction fix the original vertex and edge bounds, allow deleted labels, and induct on the number of active edges with the induction predicate conjoined over all bounded incidence tables of the given size. Equalities of
polynomials mean bounded slot equalities, connectivity is a Boolean
matrix computation, and component counts use bounded sums. In particular
the minimal-crossing conclusion quantifies over a proposed alternative
diagram and an explicit move list. It does not require producing a
shortest list or deciding whether two arbitrary diagrams are equivalent.

\paragraph{Lean proof structure.}
Use $\mathtt{Nat}$ for mask and coded input, an offset coefficient vector for
$\mathtt{Laurent}$, and a finite array of mates for a boundary pairing.
Implement \texttt{stack} by bounded component/path tracing, returning
the new pairing and a loop count. Prove its gluing law by mapping each
path segment in the two-stage trace to its original edge list. For the
local tables it suffices to tabulate $I,e,f,ef,fe$ on six endpoints
and check the three cap-cup equations by computation. For adequacy,
first establish \texttt{switch\_merge} on two different cycle IDs and
\texttt{switch\_bound} on an arbitrary pair by cutting the two edges
and tracing the two resulting paths. Define \texttt{treeMask} by the
edge scan and maintain that the current forest partition equals
the cycle partition of the tree-switched state. Count face orbits and
use the checked Euler equality $a+d=c+2$. Finally reduce a
coefficient vector into four signed slots using $A^4=-1$; the tree
bijection identifies one slot with $\pm\tau$. Each bounded invariant is
suitable for quantifier-free induction after replacing list quantifiers
by Boolean scans.

The companion \texttt{jones\_states.py} directly enumerates braid states by
gluing the levels of the braid diagram. Its open-pairing checks retain
boundary information in moves II and III. Its closed computations test the
two-strand formula and both mirror choices through eleven crossings, the
figure-eight calculation, and the value and exponent grading on all 341
three-strand words of length at most four.
